\section{DBSCAN Advanced Cluster Analysis}

\medskip

Some key R commands are:
\begin{itemize}
    \item \texttt{db <- dbscan(data, eps = eps, minPts = minPts)}: the main function call. Here, \texttt{eps} (epsilon) is the fixed radius of the neighborhoods in which DBSCAN searches for dense regions, and \texttt{minPts} is the minimum number of points required within that radius (including the point itself) for a point to be considered a core point.

    \item \texttt{kNNdistplot(data, k = k)}: produces a $k$-nearest neighbor distance plot. For each point, the distance to its $k$th nearest neighbor is computed and plotted in sorted order. This is commonly used to help choose a reasonable value for \texttt{eps}. A sharp change or ``elbow'' in the plot often suggests a good cutoff for the neighborhood radius.

    \item \texttt{table(db\$cluster)}: tabulates the number of points assigned to each cluster. In the output, cluster labels are given as integers, and points labeled as \texttt{0} correspond to noise (points not assigned to any cluster).
\end{itemize}

\medskip


Density-Based Spatial Clustering of Applications with Noise (DBSCAN) is a clustering method based on the idea that clusters are regions of high point density separated by regions of low density. Unlike methods such as $k$-means, DBSCAN can explicitly classify observations as:

\begin{itemize}[label=--]
    \item Core: points whose $\epsilon$-neighborhood (including the point itself) contains at least \texttt{minPts} points.
    \item Boundary: points that are not core points but lie within the $\epsilon$-neighborhood of a core point.
    \item Noise: points that are neither core nor boundary, and therefore do not belong to any cluster.
\end{itemize}

DBSCAN does not assume any underlying probability distribution. Instead, it identifies clusters purely through geometric density, using $\epsilon$ as a distance threshold and \texttt{minPts} as a minimum local density requirement.


\bigskip

\subsection{How the Algorithm Works}

\medskip

\begin{enumerate}[itemsep=0.5em]
    \item A data point is selected and the number of points within $\epsilon$ of it is counted. If its $\epsilon$-neighborhood (including the point itself) contains at least \texttt{minPts} points, it is classified as a ``core'' point and used as a seed for a new cluster. Otherwise, it is temporarily labeled as noise.
    
    \item The cluster is expanded by adding all points within $\epsilon$ of the seed (these points are density-reachable from the seed).
    
    \item For each point added to the cluster, if it is also a core point, then all of its $\epsilon$-neighbors are added as well. This process continues until no new density-reachable points can be added.
    
    \item Once the cluster stops growing, a new unvisited point is selected and tested to determine whether it can start another cluster.
    
    \item After all points have been processed, any point not assigned to a cluster is labeled as noise.
\end{enumerate}

This allows DBSCAN to form clusters with complex shapes by chaining together regions of high density. Compared to $k$-means, DBSCAN does not penalize points that are far apart in the same cluster, as long as there is a sufficiently dense ``path'' of points connecting them.

\bigskip

\subsection{Advantages, Applications, \& Limitations}

\medskip

The key advantages of DBSCAN are the flexibility it provides in defining density requirements through \texttt{eps} and \texttt{minPts}, its ability to explicitly identify noise points, and its capacity to recover clusters with arbitrary, non-convex shapes. Unlike $k$-means, it does not require specifying the number of clusters in advance. DBSCAN has wide applications in geospatial analysis, image processing, anomaly detection, and customer segmentation.

Key limitations of DBSCAN are its sensitivity to the choice of parameters, particularly \texttt{eps}. Small changes in \texttt{eps} can substantially alter the clustering structure. In addition, DBSCAN uses a single global $\epsilon$ value, which makes it difficult to correctly identify clusters that have different densities or widths within the same dataset. In high-dimensional settings, distance measures become less informative due to the curse of dimensionality, which can make density-based neighborhood definitions unreliable. Finally, DBSCAN can struggle when clusters are very close together, as it may incorrectly merge them if their high-density regions connect.

\bigskip

\subsection{Choosing Parameter Values}

\medskip

Since the choice for the parameters significantly affects the output, we should consider some rules of thumb:
\begin{itemize}[label=--, itemsep=0.7em]
    \item \texttt{minPts} = $p + 1$, where $p$ is the number of features. 
    \begin{itemize}[label=$\circ$, itemsep=0.3em]
        \item This is a theoretical lower bound: in $p$-dimensional space, at least $p+1$ points are required to form a full-dimensional simplex (the smallest set of points that spans the entire $p$-dimensional space rather than lying in a lower-dimensional subspace). 
        \item With fewer than $p+1$ points, observations can lie on a lower-dimensional manifold and spuriously appear dense, so DBSCAN may identify unstable or degenerate clusters. Thus $p+1$ is the smallest value that meaningfully defines local density in $p$-space, but it is rarely sufficient in practice.
    \end{itemize}
    
    \item \texttt{minPts} = $2p$
    \begin{itemize}[label=$\circ$, itemsep=0.3em]
        \item This is a common heuristic that increases the minimum neighborhood size to make density estimates more stable and less sensitive to noise. However, it is not theoretically justified: the appropriate value depends on the intrinsic dimensionality, the noise level, and the actual spatial distribution of the data.
        \item In higher dimensions or noisy datasets, substantially larger values may be required.
    \end{itemize}
    \item Experiment with a range of $\epsilon$ and \texttt{minPts} (grid search)
    \item Use relevant domain knowledge for the subject
    \item Use a \textit{k-distance graph}.
\end{itemize}

\medskip

\subsubsection{k-distance Graphs}

\medskip

This is used to \textit{refine} $\epsilon$ once we already have a tentative value for \texttt{minPts}. The idea is as follows:

\begin{enumerate}
    \item For each data point, compute the distance to its $k$th nearest neighbor, where $k$ is the chosen value for \texttt{minPts}. We record only this $k$th distance, since it represents the smallest radius needed for the point to have at least \texttt{minPts} neighbors. Distances to closer points are not needed for this diagnostic.

    \item Collect all of these $k$-nearest neighbor distances and sort them from smallest to largest.
    
    \item Plot the ordered distances against their index (from $1$ to $n$). The resulting curve will be non-decreasing (monotonically increasing).
\end{enumerate}

\medskip

To choose $\epsilon$, look for a noticeable ``elbow'' in the curve. The distance value at this elbow is often a good choice for $\epsilon$. Intuitively, the elbow marks the transition between points in dense regions (small $k$-distances) and points in sparser regions or noise (rapidly increasing $k$-distances).

